\documentclass[12pt,a4paper]{report}
\usepackage[utf8]{inputenc}
\usepackage{graphicx}
\usepackage{geometry}
\usepackage{hyperref}
\usepackage{listings}
\usepackage{xcolor}
\usepackage{setspace}
\usepackage{fancyhdr}
\usepackage{titlesec}
\usepackage{caption}
\usepackage{subcaption}

\geometry{
    a4paper,
    left=1.5in,
    right=1in,
    top=1in,
    bottom=1in
}

\definecolor{codegreen}{rgb}{0,0.6,0}
\definecolor{codegray}{rgb}{0.5,0.5,0.5}
\definecolor{codepurple}{rgb}{0.58,0,0.82}
\definecolor{backcolour}{rgb}{0.95,0.95,0.92}

\lstdefinestyle{mystyle}{
    backgroundcolor=\color{backcolour},   
    commentstyle=\color{codegreen},
    keywordstyle=\color{magenta},
    numberstyle=\tiny\color{codegray},
    stringstyle=\color{codepurple},
    basicstyle=\ttfamily\footnotesize,
    breakatwhitespace=false,         
    breaklines=true,                 
    captionpos=b,                    
    keepspaces=true,                 
    numbers=left,                    
    numbersep=5pt,                  
    showspaces=false,                
    showstringspaces=false,
    showtabs=false,                  
    tabsize=2
}

\lstset{style=mystyle}
\setstretch{1.5}

\begin{document}

\begin{titlepage}
    \begin{center}
        \vspace*{1cm}
        
        \Huge
        \textbf{Financial Explainable AI Chatbot}
        
        \vspace{0.5cm}
        \LARGE
        Project Blackbook
        
        \vspace{1.5cm}
        
        \textbf{Submitted by} \\
        \vspace{0.5cm}
        Student Name \\
        University/College Name
        
        \vfill
        
        \Large
        A project presented for the degree of\\
        Bachelor of Technology
        
        \vspace{0.8cm}
        
        \includegraphics[width=0.4\textwidth]{placeholder_logo.png} % Include a placeholder for the college logo
        \vspace{0.5cm}
        
        \Large
        Department of Computer Science and Engineering\\
        University/College Name\\
        \today
        
    \end{center}
\end{titlepage}

\pagenumbering{roman}

\chapter*{Certificate}
\addcontentsline{toc}{chapter}{Certificate}
This is to certify that the project entitled \textbf{"Financial Explainable AI Chatbot"} is a bonafide work carried out by \textbf{Student Name} in partial fulfillment for the award of the degree of Bachelor of Technology in Computer Science and Engineering from \textbf{University/College Name} during the academic year 2025-2026. The project was completed under our guidance and supervision.

\vspace{3cm}

\begin{flushleft}
    \textbf{Project Guide} \hfill \textbf{Head of Department} \\
    Name of Guide \hfill Name of HOD \\
    Designation \hfill Designation
\end{flushleft}

\vspace{2cm}
\begin{center}
    \textbf{External Examiner} \\
    ....................................
\end{center}
\pagebreak

\chapter*{Declaration}
\addcontentsline{toc}{chapter}{Declaration}
I hereby declare that the project report entitled \textbf{"Financial Explainable AI Chatbot"} submitted to \textbf{University/College Name} is a record of an original work done by me under the guidance of \textbf{Name of Guide}. This project work is submitted in the partial fulfillment of the requirements for the award of the degree of Bachelor of Technology in Computer Science and Engineering. The results embodied in this thesis have not been submitted to any other University or Institute for the award of any degree or diploma.

\vspace{3cm}

\begin{flushright}
    \textbf{Student Name} \\
    Roll Number \\
    Date: \today
\end{flushright}
\pagebreak

\chapter*{Acknowledgement}
\addcontentsline{toc}{chapter}{Acknowledgement}
I would like to express my profound gratitude and deep regards to my guide \textbf{Name of Guide} for their exemplary guidance, monitoring, and constant encouragement throughout the course of this thesis. The blessing, help, and guidance given by them time to time shall carry me a long way in the journey of life on which I am about to embark.

I also take this opportunity to express a deep sense of gratitude to \textbf{Name of HOD}, Head of the Department of Computer Science and Engineering, for their cordial support, valuable information, and guidance, which helped me in completing this task through various stages.

Lastly, I thank almighty, my parents, and friends for their constant encouragement without which this assignment would not be possible.

\vspace{2cm}

\begin{flushright}
    \textbf{Student Name}
\end{flushright}
\pagebreak

\chapter*{Abstract}
\addcontentsline{toc}{chapter}{Abstract}
Artificial Intelligence has penetrated numerous sectors, fundamentally altering how operations are executed. However, in sensitive domains like finance, the "black box" nature of complex AI models creates a barrier to trust and regulatory compliance. Explainable AI (XAI) addresses this by providing human-understandable justifications for AI decisions. 

This project, titled "Financial Explainable AI Chatbot," proposes and develops an interactive system that leverages machine learning to yield financial insights while prioritizing transparency. The chatbot integrates real-time information retrieval methods, custom financial calculation tools (Simple Interest, Compound Interest, SIP, and FD computations), and an Explainable AI engine specifically dedicated to loan eligibility scoring. 

Powered by a decoupled architecture comprising a Flask backend and a Streamlit frontend, the system orchestrates user queries using intent routing and returns both computed answers and the logic behind them. Our loan engine features dual-mode functionality: predictive capability via a pre-trained machine learning model and a robust fallback mechanism utilizing a transparent, rule-based approach when models are unavailable. The user interacts through a modern natural language AI chat interface, breaking the barrier of traditional, clunky forms and rigid web portals.

This thesis details the architectural design, algorithmic implementation, and evaluation of the system, underscoring its capacity to foster user trust in financial technologies through interactive XAI.

\pagebreak

\tableofcontents
\pagebreak

\listoffigures
\addcontentsline{toc}{chapter}{List of Figures}
\pagebreak

\pagenumbering{arabic}

\chapter{Introduction}
\section{Background}
The explosion of Artificial Intelligence (AI) algorithms in finance has revolutionized aspects of trading, risk assessment, fraud detection, and customer support. Institutions utilize powerful models to assess creditworthiness in milliseconds. However, the models achieving the highest accuracy—like deep neural networks and ensemble trees—often act as "black boxes." A user denied a loan receives a rejection without knowing whether it was due to a faulty input, minor credit history issues, or systemic algorithmic bias. Explainable AI (XAI) seeks to make AI models transparent, offering justifications behind prediction outcomes that end users and regulatory bodies can understand. Incorporating this transparency into a conversational interface fundamentally shifts the user experience from confronting an opaque system to participating in an informative dialogue.

\section{Problem Statement}
Financial institutions deploy machine learning models that are frequently incomprehensible to the end users they impact. When a user requests to check their loan eligibility, a simple "approved" or "rejected" outcome causes frustration and mistrust. Furthermore, legacy financial systems usually require users to navigate complex forms for basic calculations (such as SIP returns or Compound Interest) or to check real-time stock prices. There is an evident lack of unified, intuitive interfaces where users can effortlessly obtain financial insights combined with machine learning outputs that they can trust through transparent explanations. 

\section{Objectives}
The primary objective of this project is to build an Explainable AI Chatbot tailored for the financial sector. Specific goals include:
\begin{itemize}
    \item Designing an intuitive conversational UI utilizing Streamlit.
    \item Creating a scalable, multi-layered Python API backend orchestrating predictive models and real-time APIs.
    \item Developing a transparent loan prediction engine, complete with an elegant fallback rule-based system for high-availability.
    \item Integrating a suite of financial calculators natively responsive to natural language prompts.
    \item Enabling real-time stock market data fetching into the chat conversation.
    \item Providing verbose, explainable text outputs outlining exactly why a machine learning system made a specific financial decision.
\end{itemize}

\section{Scope of the Project}
This system encompasses personal financial domains: basic investment calculations (SIP, FD, RD), real-time stock quotes, and personal loan processing evaluation. It does not act as a licensed fiduciary, and predictions are generated for educational and demonstrative purposes. The architecture emphasizes extensibility, allowing future integration of more sophisticated XAI libraries (such as SHAP or LIME) and Large Language Models (LLMs) for reasoning.

\chapter{Literature Review}
\section{Evolution of AI in Finance}
Artificial Intelligence has been applied in finance for decades, beginning with simple expert systems in the 1980s, evolving to algorithmic trading in the 2000s, and recently permeating personal finance via mobile applications. State-of-the-art models handle enormous datasets, identifying patterns impossible for human analysts. 

\section{The Need for Explainability (XAI)}
With regulatory environments like the EU's General Data Protection Regulation (GDPR) enforcing a "Right to Explanation," companies cannot reject applicants without justification. Explainable AI encompasses tools and frameworks to dissect AI reasoning. The project incorporates foundational XAI principles through transparent rule formulation and explicit reasoning outputs mapped to the user interface.

\section{Conversational Interfaces}
Chatbots have moved from rigid decision trees to highly adaptable natural language systems. Modern frontends like Streamlit provide rapid prototyping environments for conversational agents, allowing robust API connections via libraries such as `requests` and `Flask`.

\chapter{System Design and Architecture}
\section{System Overview}
The system acts as a conversational intermediary. A user submits a query which is routed by the core engine. Depending on intent, the engine interfaces with distinct sub-modules before standardizing the response and sending it to the user.

\section{Module Breakdown}
\subsection{The Streamlit Frontend (ui.py)}
Responsible for tracking chat history, capturing user input via `st.chat_input`, and rendering messages securely. It implements robust error boundaries handling potential backend outages.

\subsection{The Flask API (app.py)}
A lightweight RESTful WSGI web application handling POST requests at the `/chat` route. It validates JSON payloads and wraps the internal engine components asynchronously.

\subsection{The Routing Engine (financial\_xai/engine.py)}
Acts as the central nervous system. Using keyword heuristics (and capable of scaling to NLP intent classification), it parses the user's sentence and delegates the execution to specific handlers (Loan, Stock, or Calculations).

\subsection{Financial and XAI Modules}
\begin{itemize}
    \item \textbf{modeling.py}: Manages the machine learning loader with failover to rule-based fallback evaluations.
    \item \textbf{stock\_service.py}: Ingests Yahoo Finance (`yfinance`) data for live equity requests.
    \item \textbf{calculations.py}: Contains hardened business logic formulas for compounding, principal, and return calculations.
\end{itemize}

\chapter{Methodology and Implementation}
\section{Tech Stack Selection}
Python was selected due to its unparalleled dominance in the machine learning ecosystem. Flask ensures minimal overhead for internal APIs, while Streamlit dramatically reduces frontend developmental time, providing innate chatbot components. the `yfinance` library handles unpredictable web scraping protocols gracefully.

\section{Loan Prediction Engine}
The engine was built accommodating physical `.pkl` models to demonstrate MLOps lifecycle integration. Should the model be absent, a Fallback Loan Predictor engages. This predictor uses transparent thresholds:
\begin{itemize}
    \item \textbf{Credit Score Thresholds}: Minimum 600 required.
    \item \textbf{DTI (Debt-to-Income) Ratio}: Computed and restricted to ensure applicants are not over-leveraged.
\end{itemize}

\section{Explainable Formatting}
When returning a prediction, the response concatenates the raw outcome (Approval/Rejection) with a structured dictionary tracing the factors contributing to the decision. This mimics local model interpretability techniques, generating immediate qualitative insight.

\chapter{Results and Discussion}
\section{Performance}
The decoupled design allows horizontal scaling. Flask comfortably handles concurrent simulation tests globally via Waitress or Gunicorn (when deployed to production). 
\section{User Experience}
Streamlit provides a pristine, familiar chat cadence. Returning real-time stock details integrated directly into an analysis pipeline grants the user unified intelligence without application switching. The fallback loan predictor proved extremely resilient during testing when model artifacts were intentionally stripped.

\chapter{Conclusion and Future Scope}
\section{Conclusion}
This project successfully demonstrates the integration of multiple asynchronous data sources, rule-based engines, and predictive architectures into a single cohesive conversational UI. It acts as a comprehensive proof of concept for the transition toward transparent financial systems, showcasing that AI need not be an opaque authority, but rather an explainable, collaborative advisor.

\section{Future Work}
Future enhancements include wrapping the engine's output in an LLM reasoning layer to naturally word the heuristic and explainable outputs. Visual XAI (like SHAP capability) is scoped to be embedded back into the UI via Streamlit's innate charting features.

\pagebreak

\begin{thebibliography}{9}
\addcontentsline{toc}{chapter}{Bibliography}

\bibitem{xai} 
Arrieta, A. B., et al. (2020). 
\textit{Explainable Artificial Intelligence (XAI): Concepts, taxonomies, opportunities and challenges toward responsible AI.} 
Information Fusion, 58, 82-115.

\bibitem{streamlit}
Streamlit Open Source Data App Framework.
\textit{Documentation and Guides}. Retrieved from \url{https://docs.streamlit.io}

\bibitem{flask}
Pallets Projects. 
\textit{Flask Web Development}. Retrieved from \url{https://flask.palletsprojects.com/}

\bibitem{yfinance}
Aroussi, R. 
\textit{yfinance: Yahoo! Finance market data downloader}. Retrieved from \url{https://github.com/ranaroussi/yfinance}

\end{thebibliography}

\pagebreak

\appendix
\chapter{Source Code Extracts}
\addcontentsline{toc}{chapter}{Source Code}

This appendix provides excerpts from the central source code to validate the system architecture.

\section{Flask Entry Point (app.py)}
\begin{lstlisting}[language=Python]
from flask import Flask, request, jsonify
from financial_xai.engine import handle_message

app = Flask(__name__)

@app.route('/health', methods=['GET'])
def health_check():
    return jsonify({"status": "Healthy"})

@app.route('/chat', methods=['POST'])
def chat():
    data = request.json
    response_text = handle_message(data.get("message", ""))
    return jsonify({"response": response_text})

if __name__ == '__main__':
    app.run(debug=True, port=5000)
\end{lstlisting}

\section{XAI Loan Modeling (modeling.py)}
\begin{lstlisting}[language=Python]
class FallbackLoanPredictor:
    def predict(self, features_dict):
        score = features_dict.get('credit_score', 0)
        income = features_dict.get('monthly_income', 1)
        debt = features_dict.get('monthly_debt_payments', 0)
        
        dti = debt / income
        
        reasons = []
        approved = True
        
        if score < 650:
            approved = False
            reasons.append("Credit score is below the minimum threshold of 650.")
        else:
            reasons.append("Credit score is acceptable.")
            
        if dti > 0.4:
            approved = False
            reasons.append("Debt-to-Income ratio exceeds the 40% maximum allowable limit.")
        else:
            reasons.append("Debt-to-Income ratio is healthy.")
            
        return approved, reasons
\end{lstlisting}

\end{document}
